\DocumentMetadata{
  pdfversion=2.0,
  pdfstandard=ua-2,
  lang=en-US,
  tagging=on,
  tagging-setup={math/setup=mathml-SE,tabsorder=structure}
}
\documentclass[11pt,letterpaper]{article}
\usepackage[margin=0.68in,includefoot]{geometry}
\usepackage{newcomputermodern}
\usepackage{amsmath,amscd,mathtools}
\usepackage{tikz,adjustbox}
\providecommand{\square}{\mdlgwhtsquare}
\usepackage[hidelinks]{hyperref}
\hypersetup{
  pdftitle={Numerical Analysis PhD Exam, August 2021},
  pdfauthor={University of Florida Department of Mathematics},
  pdfsubject={Graduate mathematics examination},
  pdfdisplaydoctitle=true,
  bookmarksopen=true
}
\setcounter{secnumdepth}{0}
\setlength{\parindent}{0pt}
\setlength{\parskip}{0.28em}
\setlength{\emergencystretch}{3em}
\setlength{\leftmargini}{2.15em}
\setlength{\leftmarginii}{2.65em}
\setlength{\leftmarginiii}{3em}
\pagestyle{empty}
\raggedbottom
\allowdisplaybreaks
\NewTaggingSocketPlug{block/list/label}{examproblem}{%
  \tagstructbegin{tag=Lbl}\tagstructend%
  \tagstructbegin{tag=\LBody}%
  \tagstructbegin{tag=\ProblemHeadingTag,title-o={\ProblemHeadingText}}%
  \tagmcbegin{tag=\ProblemHeadingTag,actualtext={\ProblemHeadingText}}%
  #2%
  \tagmcend%
  \tagstructend%
}
\newenvironment{examproblems}[2]{%
  \def\ProblemHeadingTag{#2}%
  \AssignTaggingSocketPlug{block/list/label}{examproblem}%
  \begin{enumerate}%
  \setcounter{enumi}{#1}%
  \renewcommand{\labelenumi}{\textbf{\theenumi.}}%
  \setlength{\topsep}{0.35em}%
  \setlength{\partopsep}{0pt}%
  \setlength{\itemsep}{0.48em}%
  \setlength{\parsep}{0.18em}%
}{\end{enumerate}}
\newenvironment{examsubparts}{%
  \AssignTaggingSocketPlug{block/list/label}{default}%
  \begin{enumerate}%
  \setlength{\topsep}{0.18em}%
  \setlength{\partopsep}{0pt}%
  \setlength{\itemsep}{0.16em}%
  \setlength{\parsep}{0.1em}%
}{\end{enumerate}}
\newenvironment{examinstructions}{%
  \par\begingroup\itshape
}{%
  \par\endgroup\vspace{0.18em}
}
\makeatletter
\renewcommand\subsection{\@startsection{subsection}{2}{0pt}%
  {-1.25ex plus -.25ex minus -.1ex}%
  {0.55ex plus .1ex}%
  {\normalfont\large\bfseries}}
\renewcommand{\@maketitle}{%
  \newpage\null\vskip -1em%
  {\footnotesize\bfseries
    \href{https://www.ufl.edu/}{University of Florida}, \href{https://math.ufl.edu/}{Department of Mathematics}\par}%
  \vskip -0.5em%
  \tagstructbegin{tag=H1,title={Numerical Analysis PhD Exam, August 2021}}%
  \tagpdfparaOff%
  \tagmcbegin{tag=H1}%
    {\LARGE\bfseries\href{https://gma.math.ufl.edu/exams/numerical-analysis-phd/}{Numerical Analysis PhD Exam}, August 2021\par}%
  \tagmcend%
  \tagpdfparaOn%
  \tagstructend%
  \vskip 0.25em%
  \tagmcbegin{artifact}%
    \hrule height 1pt%
  \tagmcend%
  \vskip 1em%
}
\makeatother
\title{Numerical Analysis PhD Exam, August 2021}
\author{}
\date{}
\begin{document}
\maketitle
\thispagestyle{empty}
\begin{examinstructions}
Do 4 (four) of the first 5 (1–5) and 4 (four) of the last 5 problems (6–10).

\par
You may use the following degree~\(n\) orthogonal polynomials as you like.
\[
\begin{array}{c|c|c}
n & \text{Legendre} & \text{Chebyshev} \\
\hline
0 & 1 & 1 \\
1 & x & x \\
2 & 3x^2-1 & 2x^2-1 \\
3 & 5x^3-3x & 4x^3-3x \\
4 & 35x^4-30x^2+3 & 8x^4-8x^2+1
\end{array}
\]
\end{examinstructions}
\begin{examproblems}{0}{H2}
\def\ProblemHeadingText{Problem 1.}
\item
Suppose \(A\in\mathbb{C}^{m\times m}\).
\begin{examsubparts}
\item[\textnormal{(a)}]
Suppose~\(A\) is normal and upper-triangular. Show that~\(A\) is diagonal.
\item[\textnormal{(b)}]
Prove that a Schur decomposition diagonalizes~\(A\) if and only if~\(A\) is normal.
\end{examsubparts}
\def\ProblemHeadingText{Problem 2.}
\item
This problem has two parts.
\begin{examsubparts}
\item[\textnormal{(a)}]
Suppose \(A\in\mathbb{C}^{m\times n}\) with \(m\ge n\). Show that \(A^*A\) is nonsingular if and only if~\(A\) has full rank.
\item[\textnormal{(b)}]
Suppose \(A\in\mathbb{C}^{m\times m}\), and let \(a_j\) be the~\(j\)th column of~\(A\). Prove the inequality
\[
|\det A|\le \prod_{j=1}^{m}\lVert a_j\rVert_2.
\]
 (Hint: it may be useful to consider the QR decomposition of~\(A\).)
\end{examsubparts}
\def\ProblemHeadingText{Problem 3.}
\item
Define the matrices~\(A\) and~\(B\) by
\[
A=\begin{pmatrix}
1&2&3\\
1&2&3\\
1&2&3\\
1&2&3
\end{pmatrix},
\qquad
B=\begin{pmatrix}
3&0&0\\
0&4&1\\
4&0&1\\
0&3&1
\end{pmatrix}.
\]
\begin{examsubparts}
\item[\textnormal{(a)}]
Find an economy singular value decomposition of~\(A\).
\item[\textnormal{(b)}]
Find an economy QR decomposition of~\(B\).
\end{examsubparts}
\def\ProblemHeadingText{Problem 4.}
\item
Let \(\{a_1,\ldots,a_n\}\) be a linearly independent set of vectors. Consider the Gram–Schmidt and modified Gram–Schmidt algorithms for computing an orthonormal basis \(\{q_1,\ldots,q_n\}\) so that \(\operatorname{span}\{q_1,\ldots,q_k\}=\operatorname{span}\{a_1,\ldots,a_k\}\), for each \(k=1,\ldots,n\).

\par
Suppose in computing \(q_2\), an orthogonalization error is committed so that \(q_2^*q_1=\varepsilon\).
\begin{examsubparts}
\item[\textnormal{(a)}]
Use the Gram–Schmidt algorithm to compute \(v_3\) so that \(q_3=v_3/\lVert v_3\rVert\). What is \(v_3^*q_2\)?
\item[\textnormal{(b)}]
Use the modified Gram–Schmidt algorithm to compute \(v_3\) so that \(q_3=v_3/\lVert v_3\rVert\). What is \(v_3^*q_2\)?
\end{examsubparts}
\def\ProblemHeadingText{Problem 5.}
\item
Suppose \(A\in\mathbb{C}^{m\times m}\).
\begin{examsubparts}
\item[\textnormal{(a)}]
Prove Gershgorin’s disk theorem: Let \(r_i=\sum_{j=1,\,j\ne i}^{m}|a_{ij}|\). Let \(D_i\) be the disk in~\(\mathbb{C}\) with center \(a_{ii}\) and radius \(r_i\). If~\(\lambda\) is an eigenvalue of~\(A\), then \(\lambda\in\bigcup_i D_i\); in other words, \(\lambda\) is in at least one of the disks \(D_i\).
\item[\textnormal{(b)}]
Suppose~\(A\) is Hermitian positive definite. Prove that an element of~\(A\) with largest magnitude lies on the diagonal.
\end{examsubparts}
\def\ProblemHeadingText{Problem 6.}
\item
Consider the function \(f(t)\) on \([-1,1]\) given by
\[
f(t)=
\begin{cases}
1+t, & -1\le t\le 0,\\
1-t, & 0\le t\le 1.
\end{cases}
\]
 Let \(P_k\) be the space of polynomials of degree at most~\(k\).
\begin{examsubparts}
\item[\textnormal{(a)}]
Find the best uniform approximation of~\(f\) in \(P_1\) over \([-1,1]\).
\item[\textnormal{(b)}]
Find the best approximation of~\(f\) in \(P_1\) in the norm induced by the inner-product \((u,v)=\int_{-1}^{1}uv\,dt\).
\item[\textnormal{(c)}]
Find the best approximation of~\(f\) in \(P_2\) in the norm induced by the inner-product \((u,v)=\int_{-1}^{1}uv\,dt\).
\end{examsubparts}
\def\ProblemHeadingText{Problem 7.}
\item
Let \(\{\phi_k\}_{k=0}^{n+1}\) be a set of orthogonal polynomials on \([a,b]\), with respect to the inner-product \((f,g)=\int_a^b f(x)g(x)w(x)\,dx\), indexed so that \(\phi_k\) is of degree~\(k\). Recall that the weight function~\(w\) for the inner-product satisfies \(w\in L^\infty\) and \(w(x)>0\) for almost every \(x\in[a,b]\). Prove that \(\phi_k\) has~\(k\) distinct roots \(\{x_j^{(k)}\}_{j=1}^k\), with \(x_j^{(k)}\in[a,b]\), \(j=1,\ldots,k\).
\def\ProblemHeadingText{Problem 8.}
\item
Find \(x_0,x_1,x_2\in[a,b]\), and \(c_0,c_1,c_2\in\mathbb{R}\), that maximize the degree of exactness for the quadrature formula
\[
\int_a^b f(x)\,dx\approx c_0f(x_0)+c_1f(x_1)+c_2f(x_2).
\]
\def\ProblemHeadingText{Problem 9.}
\item
The Littlewood–Salem–Izumi constant \(\alpha_0\) is the unique solution \(\alpha_0\in(0,1)\) of
\[
\int_0^{3\pi/2}\frac{\cos(t)}{t^\alpha}\,dt=0.
\]
 Describe how you could use Newton’s method together with a composite 2 point Gauss quadrature rule over~\(N\) subintervals to approximate \(\alpha_0\). Be specific about how each function used in the Newton method is defined, and be specific about how the nodes and weights for the quadrature rule are defined.
\def\ProblemHeadingText{Problem 10.}
\item
Consider the ODE \(y'=f(x,y)\) on an interval \(a\le x\le b\), with \(h=(b-a)/N\) for some given value of~\(N\).
\begin{examsubparts}
\item[\textnormal{(a)}]
Suppose~\(f\) is \(C^1\) on \([a,b]\). Compute the local truncation error (also called truncation error) for the forward Euler method \(y_{k+1}=y_k+hf(x_k,y_k)\).
\item[\textnormal{(b)}]
Suppose \(f=y^\lambda\) and \(y(a)=1\). For what values of~\(\lambda\) will the forward Euler method converge for~\(h\) small enough? You are not required to state how small~\(h\) needs to be. If your solution requires any additional assumptions, they should be clearly stated.
\end{examsubparts}
\end{examproblems}
\end{document}
